\chapter{Planaridad, mapas y coloraciones}
\label{md:planaridad}

\section{Mapas y fórmula de Euler}
\label{md:euler-plano}

Un grafo es \textbf{planar} si admite un dibujo en el plano en el que
las aristas solo se encuentran en extremos comunes.
Un \textbf{mapa} es una representación concreta de ese tipo.
Un dibujo con cruces no demuestra que el grafo sea no planar:
podría existir otro sin cruces.

Las \textbf{regiones} o caras son las componentes del plano que queda
al quitar el dibujo. Se incluye siempre la región exterior.
El grado de una región cuenta las incidencias de aristas en su frontera;
si una arista tiene la misma región a ambos lados, se cuenta dos veces.
Cada arista aporta dos lados, por lo que
\[
 \sum_{\text{regiones }R}\gr(R)=2|E|.
\]

Para un mapa de un grafo conexo finito, con $v\ge1$ vértices,
$e$ aristas y $f$ regiones, la \textbf{fórmula de Euler} es
\[
 \boxed{v-e+f=2.}
\]
Su idea es eliminar aristas de ciclos: disminuyen en uno las aristas
y las regiones, sin cambiar $v-e+f$. Se llega a un árbol, con una
región; al eliminar hojas se comprueba $e=v-1$.

Para un grafo simple conexo planar con $v\ge3$,
\[
 e\le3v-6.
\]
Si además no tiene triángulos, $e\le2v-4$.
Estas cotas usan los grados mínimos de las fronteras y Euler.
Son condiciones necesarias, no criterios suficientes de planaridad.
En $K_5$, $e=10>3\cdot5-6$; en $K_{3,3}$, que no tiene triángulos,
$e=9>2\cdot6-4$. Ambos son no planares.

\section{Subdivisiones y grafos duales}

Una \textbf{subdivisión} sustituye una arista por dos, añadiendo
un nuevo vértice intermedio; se puede repetir esta operación.
El \textbf{teorema de Kuratowski} afirma que un grafo finito es planar
si y solo si no contiene un subgrafo isomorfo a una subdivisión
de $K_5$ o de $K_{3,3}$.
La obstrucción puede estar formada por caminos que sustituyen
las aristas, sin que aparezcan literalmente esos grafos completos.

El \textbf{dual} de un mapa conexo tiene un vértice por región
y una arista por cada arista del mapa original, uniendo las regiones
situadas a sus dos lados. Puede tener aristas paralelas o lazos:
si una arista no separa regiones distintas, produce un lazo en el dual.
El dual depende del mapa elegido.
La adyacencia de regiones exige compartir una arista de frontera;
un contacto únicamente en un punto no basta.

\section{Coloraciones y grafos bipartitos}
\label{md:coloracion}

Una \textbf{coloración propia con $k$ colores} asigna a cada vértice
un color de $\{1,\ldots,k\}$, de modo que vértices adyacentes
tengan colores distintos. No es obligatorio utilizar los $k$ colores.
El problema de coloración busca el menor número necesario.

El \textbf{teorema de los cuatro colores} afirma que todo grafo simple
planar admite una coloración propia con a lo sumo cuatro colores.
En un mapa se colorean las regiones de modo que las que comparten
una arista tengan distinto color. El dual traduce este problema a
vértices, considerando solo adyacencias entre regiones distintas:
los lazos no imponen un color distinto a la propia región.

Un grafo es \textbf{bipartito} si sus vértices se dividen en dos conjuntos
sin aristas internas a ninguno. Equivale a admitir una coloración
propia con dos colores y, para grafos sin lazos,
\[
 \boxed{G\text{ bipartito}\ \eqv\ G\text{ no tiene ciclos impares}.}
\]
Un ciclo impar impide alternar dos colores y cerrar correctamente.
En ausencia de ciclos impares, dentro de cada componente se pueden
separar los vértices por la paridad de las longitudes desde un vértice
de partida; una arista entre vértices de la misma clase produciría
un ciclo impar.

El \textbf{grafo bipartito completo} $K_{r,s}$ contiene todas las
$rs$ aristas entre sus dos partes, de tamaños $r,s$.
Una \textbf{clique} es un conjunto de vértices adyacentes dos a dos,
es decir, un subgrafo completo. Una clique de tamaño $t$ requiere
al menos $t$ colores, aunque esa cota no siempre determina el mínimo.
